\documentclass[11pt]{article}
\usepackage[margin=1in]{geometry}
\usepackage{amsmath, amssymb}
\usepackage[hidelinks]{hyperref}

\title{What \texttt{landscape\_spectrum} plots}
\date{\today}

\begin{document}
\maketitle

\section*{Setup}

Let $A\in\mathbb R^{n\times n}$ be the interior block of the $Q_1$ stiffness matrix on the
$N\times N$ grid ($N=64$, so $n=(N-2)^2=3844$), $b=Mf$ the consistent load, and
$u^\star=A^{-1}b$ computed by a direct solve in \texttt{float64}. Adam is run in \texttt{float32}
on
\begin{equation*}
   L(u)=\tfrac12\|Au-b\|_2^2, \qquad u_0=0,
\end{equation*}
producing iterates $u_k$. The quantity the figure decomposes is the \emph{error}
\begin{equation*}
   e_k \;=\; u_k-u^\star \in\mathbb R^n ,
   \qquad\text{so}\qquad e_0=-u^\star .
\end{equation*}

\section*{The basis}

On a uniform grid $A=K\otimes\tilde M+\tilde M\otimes K$ with $K,\tilde M$ the 1D stiffness and
mass matrices, and the 1D sine vectors diagonalise both. Hence the 2D sine modes
\begin{equation*}
   (v_{pq})_{ij}\;\propto\;\sin\!\Big(\frac{\pi p i}{N-1}\Big)\sin\!\Big(\frac{\pi q j}{N-1}\Big),
   \qquad p,q=1,\dots,N-2,
\end{equation*}
normalised to $\|v_{pq}\|_2=1$, are \emph{exact} eigenvectors of $A$,
\begin{equation*}
   A\,v_{pq}=\lambda_{pq}\,v_{pq},
   \qquad
   \lambda_{pq}=k_pm_q+m_pk_q,
\end{equation*}
with $\theta_p=\pi p h$, $h=1/(N-1)$, $k_p=h^{-1}(2-2\cos\theta_p)$ and
$m_p=\tfrac{h}{6}(4+2\cos\theta_p)$. They are orthonormal, so
\begin{equation*}
   e_k=\sum_{p,q}c_{pq}(k)\,v_{pq},
   \qquad
   c_{pq}(k)=\langle e_k,v_{pq}\rangle,
   \qquad
   \|e_k\|_2^2=\sum_{p,q}|c_{pq}(k)|^2 .
\end{equation*}
The $c_{pq}(k)$ are obtained from one orthonormal 2D DST-I of $e_k$; no eigensolver is used. The
eigenvector property is verified numerically at run time (relative residual $\sim10^{-13}$).

\section*{The plotted quantity}

The same construction is applied to the gradient-descent iterates, giving the companion figure
\texttt{landscape\_spectrum\_gd}.

Write $\hat c_{pq}(k)$ for the root-mean-square of $|c_{pq}(k)|$ over the $25$ modes nearest in
$\lambda$ (a denoising window along the mode axis only, normalised by the number of samples
actually inside it; the $\lambda$ axis is not binned). The colour is
\begin{equation*}
   \boxed{\;
   \log_{10}\frac{\hat c_{pq}(k)}{\|e_0\|_2}\;}
   \qquad\text{clipped to }[-5,\,0],
\end{equation*}
plotted against iteration $k$ (horizontal, log) and eigenvalue $\lambda_{pq}$ (vertical, log).
Each cell is one recorded iteration $\times$ one mode.

\paragraph{Why the normalisation is global, not per mode.}
The natural choice $|c_{pq}(k)|/|c_{pq}(0)|$ is undefined here: with $u_0=0$ the initial error is
$e_0=-u^\star$, and $u^\star$ inherits the $K^2$ frequencies of the source, so \textbf{only $16$ of
the $3844$ coefficients $c_{pq}(0)$ are nonzero}. Dividing by $\|e_0\|_2$ instead makes every cell
comparable and keeps Parseval available: the squared cell values in a column sum to
$\|e_k\|_2^2/\|e_0\|_2^2$, the squared relative error.

\paragraph{Reading it, and one trap.} Bright $=$ this mode carries error comparable to the
\emph{total} initial error; black $=$ at or below $10^{-5}$ of it. The colour is a \emph{per-mode}
amplitude, and the bands hold very different numbers of modes: $3511$ modes have $\lambda\ge1$ and
only $26$ have $\lambda<0.1$. \textbf{A broad dim band can therefore carry far more error than a
narrow bright one}, and the heatmap alone does not support quantitative statements about where the
error lives. The companion figure \texttt{landscape\_bands} plots the un-diluted quantity, the
fraction of $\|e_k\|_2^2$ in each band, and that is the one to quote.

\paragraph{What the two optimizers actually do.} They behave oppositely, and only the
gradient-descent panel matches the classical intuition.
\begin{itemize}
   \item \textbf{Gradient descent} is the textbook smoother: everything above $\lambda\approx0.1$
   is annihilated within the first iteration, and $100\%$ of the squared error sits in
   $\lambda<0.1$ from iteration $1$ through $10^5$, essentially undiminished. It ends at $46.7\%$
   relative error.
   \item \textbf{Adam} does the reverse. The smooth band dominates until $k\approx10^3$, then
   crosses over sharply around $k\approx3$--$7\times10^3$, and by $k\approx10^4$ the modes with
   $\lambda\ge1$ carry $100\%$ of the error, with the smooth coefficients down at $\sim5\times10^{-5}$.
   Adam \emph{solves} the smooth content and is left rattling in the stiff modes. This is
   consistent with its per-coordinate normalisation: it advances every coordinate by
   $\sim\mathrm{lr}$ regardless of curvature, so it has no difficulty with flat directions, and no
   way to settle in steep ones.
\end{itemize}
So the ill-conditioning of $\tfrac12\|Au-b\|^2$ defeats both, but for different reasons: GD cannot
move along the flat directions, while Adam cannot stop moving along the steep ones.

\section*{Two rendering notes}

\begin{itemize}
   \item \textbf{Cell widths are unequal by construction.} Iterations are recorded log-spaced but
   must be integers, so the first ones are $1,2,3,\dots$; on a log axis the step $1\!\to\!2$ spans
   $0.30$ decades while $10^4\!\to\!10^5$ spans $1$. The wide left-hand cells are single
   iterations, and the strong colour changes between them are real: the amplitudes genuinely
   swing by an order of magnitude per step during the first few iterations.
   \item Cell edges are the geometric midpoints between samples, with the first edge pinned at
   $k=1$. (With \texttt{shading="nearest"} the first cell is centred on $k=1$ and extends half a
   log-spacing to its left, which made the axis display iterations $<10^0$.)
\end{itemize}

\end{document}
